%%%PAGE 41 rot=0 legibility=good summary=运动关联(关联速度)：绳/杆端点速度分解多例(船与人、两物体绕滑轮、定滑轮吊物、梯子靠墙、杆与方块同轴转动)，红字“沿绳方向速度相等”“同一转轴ω相同”；页末一则动滑轮合力做功
%%%BLOCK id=041-01 ch=04 sec=关联速度 kind=method alt=02 cont=none y=0.03-0.25
\textbf{运动关联}% 页面左上角标题(带下划线框线)；页左侧另有一条红色大括号，自本块“人 V”图起括至“同轴同ω”图止

\noindent $V_{\text{船}}\cos\alpha=V_{\text{人}}$\quad 船\unclear{比}人快% 页面右上角两行(第二行中间一字辨认不清)

\begin{itemize}
\item 人匀速船加速
\item 人：减速船才有可能匀速
\end{itemize}

%FIG p041_a 0.07 0.125 0.30 0.208
\notefig[0.28]{figures/p041_a.png}
% 图内标注：人 $V$(岸上，水平绳段)；绳方向分速度 $V_{\parallel}$ 改变绳长；船速 $V_1$、夹角 $\alpha$；垂直分速度 $V_{\perp}$(斜向下箭头)
\noindent $V_{\perp}$ 绳子摆动(圆周运动)% 图正下方的标注，因紧邻右侧文字未含在图内

\noindent $V=V_1\cos\alpha$\qquad 只能分解实际速度

%FIG p041_b 0.63 0.15 0.705 0.245
\notefig[0.25]{figures/p041_b.png}
% 图内标注：船受力图——$F$(沿绳斜向上，带虚线分量)、$mg$ 向下
\noindent $F\sin\alpha=mg$
%%%END
%%%BLOCK id=041-02 ch=04 sec=关联速度 kind=method alt=none cont=none y=0.25-0.34
%FIG p041_c 0.115 0.25 0.36 0.315
\notefig[0.28]{figures/p041_c.png}
% 图内标注：绳绕过上方定滑轮连接左右两物体；左物体速度 $V_1$(水平向右)，绳与水平夹角 $\alpha$；右物体速度 $V_2$(水平向右)，绳与水平夹角 $\beta$；虚线箭头、沿绳/垂直绳的分量箭头
\[
V_1\cos\alpha=V_2\cos\beta
\]
%%%END
%%%BLOCK id=041-03 ch=04 sec=关联速度 kind=method alt=03 cont=none y=0.24-0.35
%FIG p041_d 0.42 0.235 0.605 0.335
\notefig[0.25]{figures/p041_d.png}
% 图内标注：水平面上物体 $B$ 向左运动(标 $V\cos\alpha$ 的分量)，绳绕过上方定滑轮另一端竖直挂物体 $A$(向上箭头)，滑轮处标 $V_A$，夹角 $\alpha$
\noindent $V_A=V\cos\alpha$\quad $\alpha\downarrow$\quad $\cos\alpha\uparrow$

\noindent $V_A\uparrow$\quad $\therefore T>mg$
%%%END
%%%BLOCK id=041-04 ch=04 sec=关联速度 kind=method alt=none cont=none y=0.35-0.45
%FIG p041_e 0.12 0.355 0.285 0.445
\notefig[0.25]{figures/p041_e.png}
% 图内标注：直角三角形(斜边为杆/绳)，上端点 1 处标 $V_1$ 竖直向下及夹角 $\theta$，下端点 2 处标 $V_2$ 水平向右、夹角 $\theta$、虚线箭头
\[
V_1\cos\theta=V_2\sin\theta
\]
%%%END
%%%BLOCK id=041-05 ch=04 sec=关联速度 kind=note alt=none cont=none y=0.45-0.49
\red{\struck{\unclear{难}}\ 沿绳方向速度相等}% 红笔；最左一字被红笔大叉划去(像“难”)；此红字是对上方各例的总结
%%%END
%%%BLOCK id=041-06 ch=04 sec=关联速度 kind=method alt=03 cont=none y=0.46-0.61
%FIG p041_f 0.11 0.46 0.30 0.61
\notefig[0.25]{figures/p041_f.png}
% 图内标注：定滑轮(顶部横梁)，绳一端竖直挂物体 2(向上箭头)，另一端斜拉竖直轨道内的物体 1；物体 1 处标夹角 $\theta$，粗箭头向下标 $V_1$
\noindent $V_1\cos\theta=V_2$

\noindent $\theta\downarrow$\quad $V_1\cos\theta\uparrow$\quad $\therefore V_2\uparrow$\quad 2向上加速

\noindent $T>m_2g$
%%%END
%%%BLOCK id=041-07 ch=04 sec=关联速度 kind=method alt=none cont=none y=0.60-0.70
%FIG p041_g 0.07 0.605 0.268 0.685
\notefig[0.27]{figures/p041_g.png}
% 图内标注：斜杆(长 $L$，下端 $B$ 铰在斜面/地面左下，与水平夹角 $\theta$)，杆上端点 $A$；右下方块(高 $h$)以速度 $V$ 向右，其上角与杆接触，红箭头为垂直杆的分速度；红笔沿杆描出
\noindent \fcolorbox{red!75!black}{white}{\red{同轴同$\omega$}}% 红笔圈出，由 $A$ 处括号引出

\[
\left\{\begin{aligned}
&V\sin\theta=\omega\,\dfrac{h}{\text{\struck{\unclear{$\tan\theta$}}}\,\sin\theta}\\
&V_A=\omega L
\end{aligned}\right.
\]% 分母中最左一项被涂划(像 tanθ)，其后为 sinθ；原稿此处 ω 写得像 W

\noindent \red{同一转轴 $\omega$ 相同}
%%%END
%%%BLOCK id=041-08 ch=04 sec=关联速度 kind=method alt=none cont=none y=0.71-0.82
%FIG p041_h 0.07 0.715 0.245 0.81
\notefig[0.25]{figures/p041_h.png}
% 图内标注：一个五边形物体速度 $V_2$ 由绳(与水平夹角 $\alpha$)拉着；绳另一端为人，人速度 $V_1$ 水平向右(红箭头为沿绳/垂直绳分量)，虚线沿绳延长
\[
V_2=V_1\cos\alpha
\]
%%%END
%%%BLOCK id=041-09 ch=06 sec=功·合力做功 kind=derivation alt=02 cont=none y=0.83-0.93
%FIG p041_i 0.07 0.84 0.24 0.915
\notefig[0.25]{figures/p041_i.png}
% 图内标注：左侧方块，经绳连到小圆(滑轮/动点)，力 $F$ 斜向上，两侧夹角各标 $\theta/2$，虚线延长，竖直线为边界
\[
F_{\text{合}}=2F\cos\frac{\theta}{2}\qquad
W=2F\cos\frac{\theta}{2}\cdot S\cdot\cos\frac{\theta}{2}=\red{\underline{2FS\cos^{2}\frac{\theta}{2}}}=\red{\underline{FS(1+\sin\theta)}}
\]% 后两项下方各有红色下划线；CHECK: 2cos^2(θ/2)=1+cosθ，原稿末项写作 1+sinθ(笔误或误读)，照原样转录；“cos”后上标写作“2θ”分母“2”，按 cos^2(θ/2) 转录
%%%END
%%%PAGEEND 41
